\documentclass[11pt]{article}
\usepackage[margin=1in]{geometry}
\usepackage{amsmath,amssymb}
\usepackage[colorlinks=true,urlcolor=blue,linkcolor=black]{hyperref}
\setlength{\parskip}{.45em}
\setlength{\parindent}{0pt}
\setlength{\emergencystretch}{3em}
\providecommand{\tightlist}{}
\begin{document}
\section*{Three blocks: every deadline under equal singleton
marginals}\label{three-blocks-every-deadline-under-equal-singleton-marginals}

\textbf{Theorem.} For \(n=3s\), suppose every coupon has inclusion
probability \(1/3\) and every package has at most \(s\) coupons. At
every integer deadline, the probability of complete coverage is bounded
above by that for three equally likely disjoint \(s\)-blocks. This is
the balanced \(q=3\) case of
\href{https://hexagonmath.org/pdf/2611.00814v1}{Conjecture 8.1}. It does
not require independence within a draw. In fact the proof permits
successive package laws chosen adaptively, if the conditional inclusion
probability of \emph{each} coupon at each round remains \(1/3\).

We prove a more general two-block-residual bound. Suppose every coupon
is drawn with conditional marginal \(p\), with \(0<p\le1/2\). Put
\(a=1-p\), \(b=1-2p\). Given any nonempty collection of outstanding
coupons of normalized size \(x=(\text{number outstanding})/s\le2\), the
probability that some remains after \(t\) draws is at least \[
 L_t(x)=
 \begin{cases}
 a^t,&0<x\le1,\\
 xa^t-(x-1)b^t,&1<x\le2.
 \end{cases}
\] Use \(L_t(0)=0\) and \(0^0=1\).

For \(0<x\le1\), follow any single outstanding coupon: by the
conditional-marginal hypothesis its probability of remaining outstanding
is \(a^t\). For \(1<x\le2\), proceed by induction on \(t\). At \(t=0\)
the bound is exact. The normalized outstanding count \(Y\) after one
draw obeys \(Y\ge x-1>0\), \(Y\le2\), and \(\mathbb E Y=ax\). The
induction hypothesis and convexity of the positive part give a lower
bound \[
 a^t+(a^t-b^t)\mathbb E(Y-1)_+
 \ \ge\ a^t+(a^t-b^t)(ax-1)_+.
\] If \(ax\le1\), this is at least \(xa^{t+1}-(x-1)b^{t+1}\). If
\(ax\ge1\), the difference from this target equals \(p(2-x)b^t\ge0\).
This proves the induction including the case \(b=0\).

In the three-block setting the mean package size equals the
deterministic upper bound \(s\), so every draw has size \(s\). After the
first draw exactly \(2s\) coupons are outstanding, hence \(x=2\). For
\(m\ge1\), the failure probability is at least \[
 2(2/3)^{m-1}-(1/3)^{m-1}
 =3(2/3)^m-3(1/3)^m,
\] the exact three-block failure probability by inclusion-exclusion. At
\(m=0\) both laws fail with certainty. No enumeration or numerical
approximation is required.

\textbf{Strictness and a stability term.} For iid balanced draws the
partition law is the \emph{only} maximizer at any fixed \(m\ge3\) up to
labels. In the induction, replacing the positive part by \(ax-1\) (even
when negative) shows that for \(1<x<2\) and \(t\ge1\) the actual lower
bound exceeds \(L_t(x)\) by at least \(p(2-x)b^{t-1}\). After two draws
\(A,B\) in the three-block problem, the outstanding mass is
\(Y=1+|A\cap B|/s\). On \(1\le Y\le2\), the function \(L_{m-2}\) is
affine, and \(\mathbb E Y=4/3\). Thus its expectation is exactly the
partition failure probability; this is the equality case \(x=2\) in the
induction (here that affine formula even holds at \(Y=1\)). Keeping the
extra term gives the explicit separation \[
 \Pr_\mu(T>m)-\Pr_D(T>m)\ \ge\
 3^{-(m-2)}\mathbb E\!\left[
     \left(1-\frac{|A\cap B|}{s}\right)
       \mathbf1_{\{0<|A\cap B|<s\}}\right].
\] Zero gap requires that any two supported packages be identical or
disjoint. All coupons have positive marginals, so those packages form a
partition; their common block weights must be \(1/3\). Conversely the
partition attains the bound.
\end{document}
